\title{Parameter-Efficient Orthogonal Finetuning via Factorization}

\begin{abstract}
The proliferation of large foundation models is transforming numerous domains, yet constructing these models from the ground up is often prohibitively resource-intensive. Consequently, effective strategies for adapting pretrained models to specific downstream applications have become essential. This work investigates a systematic approach to finetuning, known as Orthogonal Finetuning (OFT), which facilitates such adaptation. While OFT offers impressive generalization, it necessitates tuning a considerable number of parameters because orthogonal matrices possess inherently high dimensionality. To mitigate this inefficiency, we begin with an analysis of OFT grounded in the lens of information transmission, which leads us to formulate several critical requirements for achieving improved parameter efficiency. Drawing inspiration from the Cooley-Tukey algorithm for the fast Fourier transform—renowned for its effectiveness in information processing—we introduce a compact orthogonal parameterization leveraging butterfly architectures. This refined structure is integrated into the OFT framework, giving rise to a novel, parameter-efficient adaptation technique termed Orthogonal Butterfly (BOFT). By including OFT as a limiting instance, BOFT generalizes the traditional orthogonal finetuning paradigm. To validate our approach, we conduct comprehensive experiments on the adaptation of large vision transformers, extensive language models, and text-to-image diffusion models to a suite of downstream vision and language tasks.
\end{abstract}

\section{Introduction}

Notable recent advances, exemplified by ChatGPT~\cite{brown2020language,chen2021evaluating} and Stable Diffusion~\cite{rombach2022high}, showcase the exceptional generalization performance of large-scale foundation models. The surge in their capabilities has come alongside a steep rise in parameter counts

(\emph{e.g.}, GPT-3 consists of roughly 175B parameters), making it increasingly difficult for practitioners to develop such models from scratch. Thus, the research community is motivated to find ways to harness the power of pretrained foundation models for new tasks without the impracticality of retraining. In other words, methodologies for efficiently adapting these existing models to downstream problems are of critical importance. Three principal strategies for efficient adaptation exist: \emph{model finetuning}~\cite{hulora2022,zaken2022bitfit,guo2020parameter,raffel2020exploring,qiu2023controlling,zhang2022adaptive,chavan2023one}, which modifies only a subset of parameters in the unchanged model architecture; \emph{adapter tuning}~\cite{rebuffi2017learning,houlsby2019parameter,pfeiffer2020adapterfusion,lin2020exploring,he2021towards}, which incorporates additional trainable layers or modules into the model; and \emph{prompt tuning}~\cite{li2021prefix,lester2021power}, wherein the adaptation occurs through learnable prefix tokens appended to the input sequence. Among these techniques, model finetuning stands out due to its conceptual simplicity, robust performance, and the fact that it incurs no extra computational cost during inference.

The key concept underlying model finetuning is the preservation of similarity between the pretrained and finetuned models according to specific metrics, thereby retaining the knowledge acquired during pretraining. Common finetuning practices, such as employing a low learning rate, are motivated by the intuition that limiting parameter updates will minimize deviations between the original and updated models. Given that weight matrices exhibit intrinsic structure, a more rigorous strategy seeks to maintain their relational properties, particularly the pairwise angles among neurons. This perspective underpins the Orthogonal Finetuning~(OFT)~\cite{qiu2023controlling} framework, in which a shared orthogonal transformation is applied across all neurons within a layer, ensuring the invariance of pairwise neuron angles throughout finetuning. While OFT has shown significant potential in enhancing generalization and optimization when applied to text-to-image diffusion model finetuning~\cite{qiu2023controlling}, the high dimensionality of orthogonal matrices results in a substantial number of trainable parameters. To mitigate this, OFT utilizes a block-diagonal orthogonal structure, effectively lowering the parameter count. Nonetheless, this design brings certain drawbacks: the imposed block-diagonal sparsity fixes a specific parameter distribution, and each block applies the orthogonal transformation independently. Although this approach boosts parameter efficiency, it introduces certain inductive biases—such as reduced model expressiveness and limited ability to capture traditional linear transformations—and, crucially, there is still no principled method to determine an optimal sparsity configuration.

A central challenge in this context is constructing a dense orthogonal matrix in a way that remains parameter-efficient. At first, this might appear impractical, as forming a dense orthogonal matrix of dimension $d$ typically entails $\mathcal{O}(d^2)$ parameters. Our method diverges from the conventional approach by composing a dense orthogonal matrix from a sequence of factorized sparse matrices. This technique is motivated by the realization that parameter count can be minimized by allowing for greater computational complexity, essentially trading space for computation time. Since the orthogonal matrix is formulated through repeated multiplication of sparse matrices, each training iteration necessitates performing these multiplications multiple times. To better conceptualize the factorization, we frame the generation of a dense orthogonal matrix as a process of transmitting information. Viewed this way, constructing a dense orthogonal matrix from sparse components corresponds to propagating information across a graph structured as a grid. This information transmission analogy informs numerous strategies for sparse matrix factorization that control the parameter count while maintaining sufficient flexibility to yield dense matrices.

In seeking parameter-efficient solutions within this information transmission paradigm, we are inspired by the butterfly architecture central to the Cooley-Tukey fast Fourier transform algorithm~\cite{cooley1965algorithm}, known for its ability to efficiently facilitate information exchange between nodes~\cite{klasing1994broadcasting}. The butterfly graph in the Cooley-Tukey algorithm naturally suggests a sparse factorization strategy known as butterfly factorization. Using this approach, a $d$-dimensional dense matrix can be represented as the product of $\mathcal{O}(\log d)$ sparse matrices, with each matrix containing only $\mathcal{O}(d)$ nonzero elements. By constructing each sparse matrix to be orthogonal, we attain an orthogonal parameterization requiring just $\mathcal{O}(d\log d)$ parameters, marking a substantial improvement in parameter efficiency relative to the standard method. Utilizing this streamlined orthogonal parameterization, we introduce Orthogonal Butterfly (BOFT), a novel and parameter-efficient finetuning technique. BOFT generalizes OFT as a special instance and establishes a flexible orthogonal finetuning scheme. Notably, both block-diagonal and butterfly structures rely on sparsity to constrain the number of trainable parameters in orthogonal matrices. It is worthwhile to compare this sparsity-driven methodology to the low-rank structures used in LoRA~\cite{hulora2022}; sparse and low-rank matrices each represent fundamental categories of structured matrices~\cite{candes2011robust} for achieving parameter efficiency.

In contrast to the block-diagonal configuration utilized by OFT, which balances expressivity and regularity, BOFT leverages the butterfly structure to achieve a more continuous interpolation between matrices in the full orthogonal group (for enhanced expressivity) and the identity matrix (for increased regularity). This property allows for the selection of a more compact hypothesis space within the orthogonal group that is well-suited to downstream applications. Given that the butterfly structure forms the backbone of various efficient linear transforms—such as the discrete Fourier and cosine transforms—we posit that this structured parameterization embeds an inherent inductive bias, thereby promoting improved generalization and mitigating overfitting risks. This perspective is grounded in the observation that the butterfly structure is uniquely capable of realizing a diverse range of classic linear transforms, whereas OFT's block-diagonal structure lacks this expressivity. The main contributions of this work are summarized below:

\begin{itemize}[leftmargin=*,nosep]
    \item We tackle parameter efficiency in orthogonal finetuning by introducing a novel information transmission framework, which reframes the challenge of designing a dense, parameter-efficient orthogonal matrix as an information flow problem on a grid-structured graph.
    \item Building on the butterfly structures underlying the Cooley-Tukey algorithm, we present Orthogonal Butterfly—a parameter-efficient method for orthogonal finetuning—within the context of this information transmission framework.
    \item We furnish several theoretical explanations that clarify how BOFT is capable of substantially decreasing the number of learnable parameters while maintaining robust expressivity and stable training. The matrix factorization intrinsic to BOFT additionally provides a compelling weight interpolation property (see Figure~\ref{fig:boft_interpolation}).
    \item For the first time, we extend orthogonal finetuning to a broad set of tasks beyond controllable text-to-image generation~\cite{qiu2023controlling}, demonstrating its effectiveness as a universal finetuning approach. Specifically, we evaluate BOFT across adaptation scenarios spanning computer vision and natural language processing. BOFT achieves superior results over leading methods, corroborating its strong parameter efficiency and generalization capacity.
\end{itemize}

\textbf{Parameter-efficient finetuning (PEFT)}. As foundation models (\emph{e.g.}, \cite{brown2020language,radford2021learning,rombach2022high,kirillov2023segment}) continue to scale in size and capabilities, the challenge of adapting these models for specific downstream tasks in a manner that conserves parameters has become a focal point of research~\cite{treviso2023efficient,ding2023parameter,lialin2023scaling}. There exists a broad spectrum of PEFT strategies~\cite{houlsby2019parameter,aghajanyan2020intrinsic,hulora2022,edalati2022krona,wang2022adamix,gheini2021cross,zaken2022bitfit,guo2020parameter,sung2021training,ansell2022composable,lester2021power,li2021prefix,vu2022spot,he2021towards,mao2021unipelt,karimi2021compacter,liu2022few,sung2022lst,chen2022parameter,jia2022visual,chen2022adaptformer,zhang2022neural,jie2023fact,lian2022scaling,luo2023towards}, among which reparameterization-based techniques~\cite{aghajanyan2020intrinsic,hulora2022,edalati2022krona} are most pertinent to our study. LoRA~\cite{hulora2022}, in particular, modifies pretrained weights by incorporating the product of two learnable low-rank matrices, yielding strong results for NLP tasks. Notably, LoRA utilizes a fixed rank for its added matrices; however, several other works \cite{zhang2022adaptive,valipour2022dylora,zhang2023increlora} explore the adaptive allocation of rank across different layers to better distribute the parameter budget. Owing to their straightforward design, low-rank weight reparameterizations have achieved widespread adoption~\cite{dettmers2023qlora,zi2023delta,chavan2023one}. Building on insights about hyperspherical energy's relationship with generalization~\cite{liu2018learning,liu2021orthogonal}, \cite{qiu2023controlling} introduced orthogonal finetuning, offering a compelling alternative for tuning text-to-image diffusion models. This approach learns an orthogonal transformation within a layer, and consistently delivers improved generalization and training stability relative to LoRA. Although OFT demonstrates robust performance, it typically requires a larger set of trainable parameters than LoRA, indicating that enhancing its parameter efficiency would be advantageous. Additionally, it remains an open question whether OFT can generalize beyond text-to-image diffusion control~\cite{qiu2023controlling} to broader adaptation challenges. BOFT addresses OFT’s parameter overhead by introducing butterfly factorization, thereby expanding the practical scope of orthogonal finetuning to general adaptation scenarios.

\textbf{Butterfly structures}. The radix-2 Cooley-Tukey algorithm provides an efficient decomposition of an $N$-point discrete Fourier transform into smaller $\frac{N}{2}$-point DFTs, inducing a recursive butterfly computation pattern represented as a sequence (product) of sparse matrices, collectively termed a butterfly matrix. This structure has found application in various domains, including the parameterization of orthogonal matrices to obviate pivoting in Gaussian elimination and promote computational efficiency~\cite{parker1995random}, as well as the stabilization of recurrent neural network training~\cite{jing2017tunable}, and kernel approximation~\cite{munkhoeva2018quadrature}. Research such as \cite{dao2019learning,dao2019kaleidoscope} focuses on learning fast linear transformations with butterfly-based parametrizations, while other studies \cite{chen2022pixelated,dao2022monarch} leverage butterfly matrices to facilitate more efficient neural network training. Additional use cases include accelerated matrix-vector multiplication~\cite{michielssen1996multilevel,o2010algorithm}, compact matrix approximations~\cite{li2015butterfly}, and optimization of network transmissions~\cite{klasing1994broadcasting,li2003linear,rajasingh2016transmission}. In this work, our primary aim is to exploit butterfly structures to improve the parameter efficiency of OFT, distinguishing our approach from previous research.

\section{An Information Transmission View on Orthogonal Finetuning}

We begin by outlining key concepts in OFT. When fine-tuning a pretrained weight matrix, OFT restructures the new matrix as the product of a trainable orthogonal matrix and the fixed pretrained matrix. In contrast to LoRA, which incorporates an additive low-rank component for weight updates, OFT modifies the weights using a multiplicative orthogonal transformation. While LoRA attains parameter-efficiency via a low-rank design, the original OFT achieves this through a block-diagonal orthogonal matrix~\cite{qiu2023controlling}, wherein smaller diagonal blocks yield greater parameter savings. Figure~\ref{fig:comp_lora} offers a visual comparison between these approaches. The rationale for employing orthogonal transformations in the fine-tuning process is to maintain the angular relationships between neurons~\cite{liu2018learning,liu2021orthogonal,qiu2023controlling}, thereby largely preserving the semantic knowledge acquired during pretraining. Specifically, OFT involves optimizing an orthogonal matrix $\bm{R}\in\mathbb{R}^{d\times d}$ for a pretrained linear layer represented by $\bm{W}^0\in\mathbb{R}^{d\times n}$. The computation changes from $\bm{z}=(\bm{W}^0)^\top\bm{x}$ to $\bm{z}=(\bm{R}\bm{W}^0)^\top\bm{x}$, where $\bm{x}\in\mathbb{R}^{d}$ serves as the input vector and $\bm{z}\in\mathbb{R}^{n}$ as the output. To facilitate zero initialization, OFT sets $\bm{R}$ initially as the identity matrix. Preservation of orthogonality during the tuning phase is guaranteed by applying Cayley parameterization~\cite{qiu2023controlling,liu2021orthogonal}, specifically, $\bm{R}=(\bm{I}+\bm{Q})(\bm{I}-\bm{Q})^{-1}$, where $\bm{Q}$ denotes a skew-symmetric matrix ($\bm{Q}=-\bm{Q}^\top$). For efficiency, a block-diagonal parameterization is utilized, expressing the orthogonal matrix as $\text{diag}(\bm{R}_1,\bm{R}_2,\cdots,\bm{R}_r)$, with each block $\bm{R}_i\in\mathcal{R}^{b\times b}$ and $br=d$, ensuring each block remains orthogonal. However, this block-diagonal scheme inherently partitions each neuron’s dimensions (columns of $\bm{W}^0$) into $r$ groups, where dimensions from distinct groups are independently transformed via separate orthogonal matrices. Although such a block-wise structure performs well in practice, the index-based grouping of neuron dimensions lacks an intrinsic rationale, suggesting a preference for dense orthogonal matrices. Consequently, an important question emerges: \emph{Is it feasible to construct a dense orthogonal matrix while preserving parameter-efficiency?}

To tackle this issue, we introduce an approach where a dense orthogonal matrix is constructed as a sequence of products involving several sparse orthogonal matrices. For a cohesive and accessible analysis of sparsity structures in orthogonal matrix factorization, we reframe the task of generating a dense orthogonal matrix in OFT as analogous to an information transmission problem. Concretely, if we create a dense matrix $\bm{R}\in\mathbb{R}^{d\times d}$ by multiplying $m$ square matrices such that $\thickmuskip=2mu \medmuskip=2mu \bm{R}=\bm{B}_m\bm{B}_{m-1}\cdots\bm{B}_1$, this process can be visualized as the transfer of information across a grid of $d\times (m+1)$ nodes (see Figure~\ref{fig:info_trans}). The rationale for this viewpoint lies in recognizing that a dense $d\times d$ matrix signifies complete connectivity between two distinct sets of $d$ nodes each. In this context, a zero entry $R_{ij}$ in $\bm{R}$ denotes the absence of a transmission path from node $j$ to node $i$, while a non-zero $R_{ij}$ signifies feasible information flow. Accordingly, expressing $\bm{R}$ as the product $\bm{B}_m\bm{B}_{m-1}\cdots\bm{B}_1$ can be interpreted as a multi-stage process of information propagation, dictated by the graph structures induced by each matrix $\bm{B}_i$. Here, information flows sequentially from the connectivity dictated by $\bm{B}_1$ through to $\bm{B}_n$. In the example provided in Figure~\ref{fig:info_trans}, the specific factorization $\bm{R}=\bm{B}_5\bm{B}_4\bm{B}_3\bm{B}_2\bm{B}_1$ is illustrated, including its sparsity pattern and corresponding graphical representation. The depicted graph arises from a step-by-step unrolling of matrix multiplication; within this framework, each $\bm{B}_i$ acts as an adjacency matrix defining connections from nodes at the $i$-th level to those at the $(i+1)$-th level. Explicitly, the $(j_1,j_2)$ entry in $\bm{B}_i$ indicates whether there exists a directed edge from the $j_2$-th node at level $i$ to the $j_1$-th node at level $(i+1)$ (with a zero indicating no such connection). For the overall product $\bm{B}_5\bm{B}_4\bm{B}_3\bm{B}_2\bm{B}_1$ to result in a dense matrix, every node at the initial level must be able to convey information to all nodes at the sixth level. If instead we analyze only the partial product $\bm{R}=\bm{B}_4\bm{B}_3\bm{B}_2\bm{B}_1$, which links source nodes at the first level to recipient nodes at the fifth level, it is evident that, for instance, node 1 is unable to transmit information to node 3. Consequently, this is reflected in the sparsity structure of $\bm{B}_4\bm{B}_3\bm{B}_2\bm{B}_1$, where the entry at position $(3,1)$ is zero. The same correspondence between induced graphical connections and the sparsity pattern holds for shorter products, such as $\bm{R}=\bm{B}_3\bm{B}_2\bm{B}_1$ or $\bm{R}=\bm{B}_2\bm{B}_1$. Broadly, for a matrix $\bm{R}\in\mathbb{R}^{d\times d}$ to be fully dense, the $m$ factor matrices $\bm{B}_m,\ldots,\bm{B}_1$ must represent directed edges within a $d\times (m+1)$ grid, ensuring that every node from the initial layer is able to transmit information to every node at the final, $(m+1)$-th, layer, with each directed edge only bridging adjacent layers (that is, connecting nodes between neighboring columns).

The perspective of information transmission offers valuable insights into the sparsity characteristics of factorization matrices within OFT. In Figure~\ref{fig:blk_diag}, the block-diagonal arrangement of $\bm{R}$ as seen in the canonical OFT is illustrated. While such a structure leads to fewer trainable parameters, it inherently restricts $\bm{R}$ from achieving full density. The central objective is to synthesize a dense orthogonal matrix from $m$ individual sparse orthogonal matrices, all while minimizing the quantity of essential trainable parameters. Adopting the information transmission viewpoint, we articulate two primary criteria to steer this objective: \emph{(i)} \textbf{dense connectivity}, which stipulates that every node at the input stage must have at least a single connecting route to each node at the output stage, and \emph{(ii)} \textbf{minimum free edges}, requiring that, under orthogonality constraints, the sum of edges remains as limited as feasible. Importantly, orthogonality imposes nuanced restrictions on the permitted edges between successive levels. Specifically, for a matrix $\bm{B}_i$ to be of full rank (a prerequisite for orthogonality), $d$ edges are necessary to establish a bijective mapping between the nodes of the $i$-th and $(i+1)$-th levels. Consequently, at least $d$ edges must connect each pair of adjacent levels (for instance, 4 in the setting of Figure~\ref{fig:info_trans}). These edges, required for orthogonality, are considered non-trainable and should thus be omitted from the trainable edge count, since such entries in a $d \times d$ orthogonal matrix are fixed at $\pm 1$ and not tunable. Given that orthogonal matrices demand fewer independent parameters compared to general dense matrices, orthogonality induces complex interdependencies among edge connections. For example, eliminating the entry at $(1,1)$ in a $2 \times 2$ orthogonal matrix necessarily removes the $(2,2)$ entry as well; in other words, omitting one edge may force the removal of another. Accordingly, certain admissible connection patterns may gain or lose edges solely due to orthogonality requirements. Employing this visualization of the sparsity pattern within the synthesized orthogonal matrix, the information transmission approach is especially advantageous for analyzing OFT, as the focus lies exclusively on the positions of non-zero trainable elements in $\bm{R}$, with their explicit magnitudes being irrelevant.

A straightforward approach to establishing dense connectivity between two hierarchical levels involves $\mathcal{O}(d^2)$ connections (i.e., utilizing a single dense orthogonal matrix), which amounts to $d^2-d$ links; for instance, when $d=4$, this results in $12$ connections. As depicted in Figure~\ref{fig:info_trans}, a possible matrix factorization is illustrated that requires only $10$ connections overall—fewer than are needed for a full dense orthogonal matrix. This approach provides a graphical lens through which to examine the parameter-efficiency of OFT and offers a systematic way to generate valid factorizations. Our method takes cues from the butterfly graph structure inherent to the Cooley-Tukey algorithm~\cite{cooley1965algorithm}, which connects $d$ sources to $d$ sinks with just $\mathcal{O}(d\log d)$ edges, achieving dense interconnection in an efficient manner. For instance, the arrangement in Figure~\ref{fig:info_trans} necessitates $10$ edges for dense connectivity, whereas a butterfly network configuration would require only $8$. The following section describes how adopting a butterfly architecture enhances parameter efficiency.

\section{Orthogonal Parameterization by Butterfly Factorization}

Originally developed as part of the Cooley-Tukey algorithm for fast Fourier transform computation, the butterfly network is instrumental in rapidly mixing information: in Fourier analysis, small changes in frequency can influence all points in the spatial domain. This principle aligns with our goal of allowing each node at the input level to communicate with every node at the output level. Moreover, the butterfly pattern is well established in computer networking~\cite{solihin2015fundamentals,li2011linear} for its efficacy in transmitting information. Let $k \geq 2$ be a power of $2$; the \emph{butterfly factor} $\bm{B}^F(k)$ is then defined as

\begin{equation}\label{eq:but_factor}
\begin{aligned}
    \bm{B}^F(k)=\begin{bmatrix}
    \text{diag}(\bm{d}_1) & \text{diag}(\bm{d}_2) \\
    \text{diag}(\bm{d}_3) & \text{diag}(\bm{d}_4)
    \end{bmatrix}\in\mathbb{R}^{k\times k},
\end{aligned}
\end{equation}

with vectors $\bm{d}_i\in\mathbb{R}^{\frac{k}{2}}, \forall i$. For $d=2^N$, the recursive definition of a $d$-dimensional \emph{butterfly matrix} $\bm{B}(d)\in\mathbb{R}^{d\times d}$ is as follows:

\begin{equation}
\begin{aligned}
    \!\!\bm{B}(d)=\tilde{\bm{B}}(d,d)\!\cdot\!\begin{bmatrix}
    \bm{B}_1(\frac{d}{2})\!\!\!\!\! & \bm{0} \\
    \bm{0}\!\!\!\!\! &\bm{B}_2(\frac{d}{2})
    \end{bmatrix}= \tilde{\bm{B}}(d,d)\tilde{\bm{B}}(d,\frac{d}{2})\cdots\tilde{\bm{B}}(d,2),
\end{aligned}
\end{equation}

Here, $\bm{B}_1(\frac{d}{2})$ and $\bm{B}_2(\frac{d}{2})$ represent butterfly matrices of dimension $\frac{d}{2}$. We refer to the \emph{butterfly component} as $\thickmuskip=2mu \medmuskip=2mu \tilde{\bm{B}}(d,k)=\text{diag}(\bm{B}^F_1(k),\cdots,\bm{B}^F_{d/k}(k))$, which forms a $d\times d$ block-diagonal matrix composed of blocks of size $k$. Here, each $\bm{B}^F_i(k)$ (for all $i$) corresponds to butterfly factors as specified in Equation~\ref{eq:but_factor}. With these definitions, we can use butterfly matrices to represent orthogonal matrices. To accomplish this, it suffices that each multiplicative factor $\tilde{\bm{B}}(d,k)$ (for all $k$) in the butterfly matrix $\bm{B}(d)$ is orthogonal. Focusing first on the case when the block size is $2$, $\tilde{\bm{B}}(d,2)$ is easily made orthogonal by using the Cayley transform (i.e., parameterizing each $2\times 2$ block via a $2$-dimensional rotation), as adopted by \cite{liu2021orthogonal, qiu2023controlling}. Since the pattern of non-zero entries in every butterfly component corresponds to a permutation of those in $\tilde{\bm{B}}(d,2)$, all such components can be parameterized as orthogonal matrices with relative ease. This construction yields a practical scheme for representing orthogonal matrices as products of several $2\times 2$ orthogonal blocks. Extending this idea as in \cite{chen2022pixelated}, we introduce a block butterfly component $\tilde{\bm{B}}^b(d,k)$ in which each element of $\bm{d}_i$ is now a $b\times b$ matrix. To ensure that $\tilde{\bm{B}}^b(d,2)$ is orthogonal, each $2b\times 2b$ block matrix is required to be orthogonal as well. Furthermore, for $k>2$, the block structure of $\tilde{\bm{B}}^b(d,k)$ arises from block-wise permutations of the pattern found in $\tilde{\bm{B}}^b(d,2)$, which also supports parameterization as orthogonal matrices. Collectively, these components form the forward computation in BOFT as follows:

\begin{equation}
    \begin{aligned}\nonumber
    \bm{z}=\big(\bm{R}(m,b)\cdot\bm{W}^0\big)^\top\bm{x},~~~~\text{s.t.}~\bigg{\{}\bm{R}(m,b)=\prod_{i=1}^m \tilde{\bm{B}}^b_{(d,i)}~~\&~~\underbrace{\big(\tilde{\bm{B}}^b_{(d,j)}\big)^\top\tilde{\bm{B}}^b_{(d,j)}=\tilde{\bm{B}}^b_{(d,j)}\big(\tilde{\bm{B}}^b_{(d,j)}\big)^\top\!=\bm{I}_d}_{\forall j\in [1, m]}\bigg{\}},
    \end{aligned}
\end{equation}

Here, for brevity, we let $\tilde{\bm{B}}^b(d,2^{m - i+1})$ be represented as $\tilde{\bm{B}}^b_{(d,i)}$, and denote by $\bm{I}_d$ the $d \times d$ identity matrix. The orthogonal matrix $\bm{R}(m,b)\in\mathbb{R}^{d\times d}$ is structured as the product of several orthogonal butterfly factors. For ease of reference, we describe BOFT using $\bm{R}(m,\frac{b}{2})$ as BOFT($m$,$b$) with the assumption $b\geq 2$. In the particular case where $m=1$, BOFT($1$,$b$) becomes equivalent to the block-diagonal OFT~\cite{qiu2023controlling} where the block size is $b$. Similarly, choosing BOFT($1$,$d$) matches the original OFT~\cite{qiu2023controlling} with a fully unconstrained orthogonal matrix. For BOFT($\log \frac{2d}{b}$,$b$), the block butterfly matrix $\bm{B}^{\frac{b}{2}}(d)$ functions as $\bm{R}$, thus yielding a dense orthogonal matrix. In a more general setting, the total number of trainable parameters in BOFT($m$,$b$) is $\frac{1}{2}(b-1)dm$ when adapting a linear layer of shape $d\times n$. If the butterfly structure is utilized (\emph{i.e.}, $m=\log d,b=2$), BOFT employs $\mathcal{O}(d\log d)$ parameters. In comparison, a standard OFT based on a full dense orthogonal matrix requires $\mathcal{O}(d^2)$ parameters, while a block-diagonal OFT with $r$ blocks needs $\mathcal{O}(bd)$. Thus, to realize a dense orthogonal matrix, the conventional OFT must set $b=d$, whereas BOFT is flexible and can accommodate arbitrary $b$ values to produce such a matrix.

\textbf{BOFT Identity Initialization}. Finetuning approaches commonly initialize with weights from the pretrained model to ensure that modifications stay close to the original configuration. For example, LoRA applies zero initialization to its low-rank adaptation matrices. In the BOFT framework, we initialize each butterfly component as the identity (which in practice means initializing the skew-symmetric matrices for the Cayley transform to zero).

\textbf{Multiplicative Dropout within BOFT}. LoRA~\cite{hulora2022} incorporates Dropout to its low-rank layers to alleviate overfitting risks. While standard Dropout~\cite{srivastava2014dropout} integrates seamlessly with LoRA, it is not suitable for BOFT due to its multiplicative parameter updates. We thus introduce a version of Dropout tailored for BOFT: the multiplicative Dropout. Since the matrix $\bm{R}(m,b)$ consists of $m$ orthogonal butterfly components, which can be rearranged as block-diagonal orthogonal matrices of size $2b \times 2b$, the multiplicative Dropout randomly selects $p_1$ proportion of the butterfly components and, within each, a $p_2$ fraction of diagonal blocks, replacing the selected blocks by identity matrices.

\section{Intriguing Insights and Discussions}

\textbf{Expressivity of BOFT}. While the butterfly structure augmented with permutations is capable of exactly replicating numerous classical fast linear transforms~\cite{dao2019learning,dao2019kaleidoscope} (\emph{e.g.}, fast Fourier transform, Hadamard transform), it remains an open question to what degree our orthogonal butterfly matrix can serve as an effective approximation for arbitrary orthogonal matrices. To investigate this, we simulate the approximation of a randomly generated dense orthogonal matrix of size $512\times 512$~\cite{anderson1987generation}. As illustrated in Figure~\ref{fig:para_eff}, results represent the mean over 10 randomly initialized trials. Here, the y-axis represents the approximation error, while the x-axis captures the effective number of tunable parameters. In each curve, color encodes BOFT variants with identical block sizes, with the leftmost data point indicating the error for BOFT($1$,$b$) (\emph{i.e.}, the standard block-diagonal OFT with block size $b$). Across all cases, BOFT generally demonstrates improved parameter efficiency relative to OFT. For instance, BOFT($9$,$2$) achieves higher expressivity compared to BOFT($1$,$16$), yet uses substantially fewer parameters. Variants with smaller $b$ and larger $m$ typically enjoy increased parameter efficiency; as an example, BOFT($6$,$4$) reaches a comparable approximation error as BOFT($2$,$16$), despite leveraging far fewer parameters. Overall, the butterfly matrix defines a more constrained and structured subclass within the orthogonal group than the block-diagonal design, which underlies why BOFT possesses provably greater expressivity than OFT with an equivalent block size.

\begin{theorem}[Expressivity of BOFT]\label{thm:exp_boft}
    BOFT surpasses OFT in expressivity when using the same block size. Any orthogonal matrix of dimension $d$ can be exactly represented by composing butterfly matrices in the form  $\bm{B}_{d-1,1}(d)\bm{B}^\top_{d-1,2}(d)\cdots\bm{B}_{1,1}(d)\bm{B}^\top_{1,2}(d)$, where each $\bm{B}_{i,j}(d)$ is a butterfly matrix for all valid $i$ and $j$.
\end{theorem}

Based on Theorem~\ref{thm:exp_boft}, we observe that BOFT can be naturally generalized—namely, the final orthogonal matrix may be written as $\thickmuskip=2mu \medmuskip=2mu \bm{R}^G(m_1,b_1,m_2,b_2,l)=\bm{R}_{l,1}(m_1,b_1)\bm{R}_{l,2}^T(m_2,b_2)\cdots\bm{R}_{1,1}(m_1,b_1)\bm{R}_{1,2}^T(m_2,b_2)$, where $\bm{R}_{i,j}^T(m,b)$ denotes the orthogonal matrix appearing in BOFT. Setting $m_1=m_2=\log d$, $b_1=b_2=2$, and $l=d-1$, ensures that $\bm{R}^G(m_1,b_1,m_2,b_2,l)$ can capture any member of the orthogonal group—an architecture referred to as the kaleidoscope hierarchy~\cite{dao2019kaleidoscope}. Nonetheless, it is important to recognize that maximizing expressivity does not guarantee optimal performance in finetuning scenarios. For example, full finetuning methods, despite being maximally expressive, frequently yield disappointing results. Thus, an optimal balance between expressivity and regularity is essential for robust model finetuning. By leveraging structural inductive biases, BOFT broadens the parameter search space during finetuning, providing a principled way to navigate the trade-off between expressivity and regularization.

\textbf{Spectral properties}. Compared to LoRA, orthogonal finetuning more effectively maintains desirable spectral characteristics, as it exactly conserves the spectral norm of the original pretrained weight matrix $\bm{W}^0$. This can be formally demonstrated via the singular value decomposition: $\bm{W}^0=\bm{U}\bm{\Sigma}\bm{V}^\top$, where $\bm{U}$ and $\bm{V}$ are orthogonal matrices, and $\bm{\Sigma}$ is diagonal with singular values. With both OFT and BOFT, an orthogonal matrix $\bm{R}$ is applied to the left, resulting in updated weights $\bm{R}\bm{U}\bm{\Sigma}\bm{V}^\top$; this operation leaves the maximal singular value—namely, the spectral norm—of $\bm{W}^0$ unchanged. This invariance has been found to enhance both stability during training and model generalization~\cite{miyato2018spectral,yoshida2017spectral}. Additional theoretical properties are discussed in Appendix~\ref{app:sn_property}.

\textbf{Orthogonal finetuning as learning bilinear similarity}. The BOFT framework can also be interpreted as optimizing a bilinear similarity function, specifically $\bm{w}^0_i\bm{R}\bm{x}$, where $\bm{w}^0_i$ denotes the $i$-th column (or neuron) of the weight matrix $\bm{W}^0$. In this context, BOFT amounts to identifying the bilinear similarity matrix $\bm{R}$ under the orthogonality constraint, which closely aligns with established concepts in distance metric learning~\cite{xing2002distance} and the theory of bilinear forms~\cite{roman2005advanced}.

\textbf{Inductive bias and generalization in BOFT}. The matrix $\bm{R}(m,b)$ employed in BOFT typically lies within a particular, structured subclass of the orthogonal group, thereby restricting the hypothesis space and introducing an inductive bias. We suggest that the particular structured inductive bias from butterfly factorization is advantageous for model generalization, since it mirrors the shared structural motifs found in a variety of well-known linear transformations~\cite{dao2019kaleidoscope}, such as the discrete Fourier transform, discrete sine/cosine transform, and Hadamard transform. Furthermore, the sparse nature of matrix factorization in BOFT may confer additional implicit inductive biases~\cite{gunasekar2017implicit,li2018algorithmic,liu2019neural}.

\textbf{Comparison to butterfly-based sparse training}. Numerous studies~\cite{dao2019learning,dao2019kaleidoscope,chen2022pixelated,dao2022monarch} have explored the use of the butterfly parameterization for sparse training. These efforts commonly involve directly expressing the weight matrices in neural networks using the butterfly structure, and conducting training processes from initialization. In contrast, \cite{dao2022monarch} investigates fine-tuning of pretrained weights by initially projecting these weights onto a variant of butterfly matrices, followed by updating the resulting components tailored for downstream applications. BOFT introduces a fundamentally distinct fine-tuning approach, where the weights undergo transformation via layer-shared weight matrices.

\input{rebuttal-results}

\section{Concluding Remarks and Limitations}

This work introduces Orthogonal Butterfly, a parameter-efficient fine-tuning framework applicable to foundation models, which leverages the butterfly structure. The central idea to enhance parameter efficiency is to express a dense orthogonal matrix as a sequence of products involving multiple sparse orthogonal matrices. To facilitate the identification of appropriate matrix decompositions, we put forward a graphical information transmission perspective. Within this framework, we observe that the butterfly pattern serves effectively to fulfill the requirements for sparse orthogonal matrix decomposition. BOFT is empirically shown to be advantageous for fine-tuning a variety of large models, including language, vision, and text-to-image generative models. The experimental outcomes substantiate BOFT's advantages as a universal solution for model fine-tuning.

Nonetheless, BOFT does have certain shortcomings. Because the orthogonal matrix produced in BOFT arises from the product of several orthogonal matrices, there is a modest increase in training computational cost compared to OFT. Addressing the challenge of reducing BOFT’s training overhead is left for future research. On the positive side, after the fine-tuning procedure, the orthogonal matrices learned by BOFT can be absorbed into the original pretrained model without introducing any additional cost during inference. Additionally, it remains an open question whether the butterfly structure is truly optimal in terms of information transfer efficiency. The proposed information transmission perspective invites insights from related fields, such as computer networking, where network topology efficiency is a well-examined topic. We anticipate that exploring alternative network structures may yield even more effective methods for constructing dense orthogonal matrices.

\end{document}